\chapter{System Architecture and Design}
\label{ch:architecture}

\section{Client-Server Architecture}

FairLint-DL operates as a client-server application with two primary components communicating over HTTP/REST on \texttt{localhost:8765}:

\begin{itemize}
    \item \textbf{VS Code Extension (Client):} Written in TypeScript, the extension handles all user interaction, file system access, WebView dashboard rendering with Plotly.js charts, and manages the lifecycle of the Python backend.
    \item \textbf{Python FastAPI Backend (Server):} Hosts the PyTorch models and all analysis algorithms, including data preprocessing, model training, QID computation, causal debugging, discriminatory instance search, and SHAP/LIME explainability.
\end{itemize}

\begin{figure}[H]
    \centering
    \includegraphics[width=0.95\textwidth]{images/fairlint-dl-architecture.png}
    \caption{FairLint-DL client-server architecture overview.}
    \label{fig:architecture}
\end{figure}

When the extension activates, it spawns the Python backend as a child process via \texttt{python -m uvicorn bias\_server:app --port 8765} and polls the health endpoint (\texttt{GET /}) until the server responds, retrying up to 15 times with 1-second intervals.

\section{REST API Endpoints}

The backend exposes seven endpoints that correspond to the six-step analysis pipeline:

\begin{table}[H]
\centering
\caption{FairLint-DL REST API endpoints.}
\label{tab:api}
\begin{tabular}{@{}llp{7cm}@{}}
\toprule
\textbf{Endpoint} & \textbf{Method} & \textbf{Description} \\
\midrule
\texttt{/} & GET & Health check \\
\texttt{/columns} & POST & Returns column names, sample data, auto-detected sensitive features \\
\texttt{/train} & POST & Trains proxy model (or loads from cache); supports \texttt{force\_retrain} \\
\texttt{/activations} & POST & Returns PCA/t-SNE reduced layer activations \\
\texttt{/analyze} & POST & Computes QID metrics and group fairness metrics \\
\texttt{/search} & POST & Runs global/local discriminatory instance search \\
\texttt{/debug} & POST & Performs layer and neuron sensitivity analysis \\
\texttt{/explain} & POST & Generates batch SHAP/LIME explanations \\
\texttt{/explain-instance} & POST & Single-instance LIME explanation (by index or custom values) \\
\bottomrule
\end{tabular}
\end{table}

\section{Project Structure}

The codebase is organized into the following directory structure:

\begin{lstlisting}[language={}, caption={FairLint-DL project structure.}]
fairlint-dl/
+-- python_backend/        # Analysis Server
|   +-- analyzers/         # Core Algorithmic Logic
|   |   +-- causal_debugger.py    # Layer/Neuron attribution
|   |   +-- qid_analyzer.py       # Entropy-based metrics
|   |   +-- search.py             # Gradient-guided search
|   |   +-- explainability.py     # SHAP and LIME explanations
|   |   +-- group_fairness.py     # DP, EO, EOpp metrics
|   |   +-- internal_space.py     # PCA/t-SNE visualization
|   +-- models/
|   |   +-- fairness_dnn.py       # PyTorch DNN definition
|   +-- utils/
|   |   +-- data_loader.py        # CSV loading, encoding, scaling
|   |   +-- model_cache.py        # SHA-256 caching system
|   +-- bias_server.py            # FastAPI entry point
+-- src/                   # VS Code Extension Source
|   +-- extension.ts              # Main entry point
|   +-- analysis/pipeline.ts      # 6-step pipeline orchestration
|   +-- webview/
|       +-- results.ts            # WebView panel management
|       +-- htmlBuilder.ts        # Dashboard HTML generation
|       +-- charts.ts             # Plotly.js chart rendering
|       +-- scoring.ts            # Composite fairness score
+-- package.json                  # Extension manifest
\end{lstlisting}

\section{Six-Step Analysis Pipeline}

The analysis follows a sequential six-step pipeline orchestrated by \texttt{pipeline.ts}:

\begin{enumerate}
    \item \textbf{Model Training:} Train (or load cached) DNN proxy model on the dataset.
    \item \textbf{Internal Space Visualization:} Extract and reduce layer activations via PCA.
    \item \textbf{QID Analysis:} Compute Shannon entropy QID, min-entropy QID, disparate impact, and group fairness metrics across all protected attributes~\cite{dice}.
    \item \textbf{Discriminatory Instance Search:} Gradient-guided global search followed by perturbation-based local search~\cite{dice, neufair}.
    \item \textbf{Causal Debugging:} Layer sensitivity analysis followed by neuron-level impact scoring~\cite{dice}.
    \item \textbf{Explainability:} Batch SHAP (KernelExplainer on log-odds) and LIME explanations with interactive per-instance explorer~\cite{kamolov2025, wang2024}.
\end{enumerate}
